\section{Discussion}
\label{sec:discussion}

\paragraph{What the nulls mean.} A reader could summarise this paper as ``foundation models mostly tie''. That
summary is accurate and is the point. The zero-shot claim is usually evidenced on aggregate leaderboards without
significance testing \citep{godahewa2021monash,aksu2024gifteval,qiu2024tfb,li2025tsfmbench}. Score the same models on a tropical station network, resample at the forecast origin, and
correct across 48 comparisons, and most margins on observations become indistinguishable from a well-built
baseline. The models are not weak: six of eight outrank every trained baseline, and 272 of 432
skill intervals exclude climatology. But the size of the advantage on real station data is smaller than the
framing in that literature implies.

\paragraph{Why rainfall is different.} The rainfall result is not a modelling failure so much as a property of the
target. Daily tropical rainfall is zero-inflated (\SI{67.6}{\percent} dry station-days) and convective, so beyond a
day the predictable signal is largely the seasonal cycle, which a day-of-year climatology already encodes. A
forecast that reproduces the climatological distribution is close to the achievable ceiling, and no model in our
roster exceeded it beyond a one-day lead. This is worth stating plainly because rainfall is the variable on which
the operational value of these models would be judged in Bangladesh, and because precipitation is difficult
enough to warrant benchmarks of its own \citep{zhang2025rainfallbench}.

\paragraph{Why regime stratification matters.} The spell result is the clearest argument against aggregate
reporting. That aggregates can hide regime-dependent failure has been shown for empirically defined regimes in
road traffic \citep{wang2026regime}; here the regimes are physically defined and externally validated
\citep{ferdoushi2023spells}, and the failure is a sign reversal rather than a degradation. The same models are strongly skilful during break spells and significantly worse than climatology during
active spells, a reversal that vanishes in a seasonal average. One mechanism would explain this, though we do not test it. A model with a
long context window carries the persistent warm, dry conditions of a break forward. That is rewarded while the
break holds, and penalised as soon as convection resumes. Whatever the cause, a benchmark that reports only aggregate skill would
certify these models as adequate for exactly the situation in which they are least reliable.

\paragraph{Implications for benchmark design.} Three of our results are really statements about how evaluation is
done. Scoring against reanalysis inflates apparent skill and, more subtly, inflates the \emph{margin} between model
classes, so reanalysis-only benchmarks \citep{rasp2024weatherbench2,nathaniel2024chaosbench} are not a neutral
convenience. \citet{gupta2026mausam} reach the same conclusion for AI weather-prediction models over South Asia;
our result transfers it to zero-shot time-series models. Stratifying by a physically defined regime
exposes failures that aggregates hide. And a context-availability control changes the interpretation of a
geographic gap: without masking the trained baseline identically, our own earlier reading, that longer context
dissolves the tropical deficit, did not survive. This matters for how a North--South divide in forecast skill
\citep{mozaffari2026inequality,masi2025safe} should be read, because part of what looks geographic is an
observing-system gap that sits inside the model's input window.

\paragraph{Affordability.} The entire zero-shot leaderboard reproduces in 4.16 GPU-hours on a single GPU, and the
strongest classical baseline we ran costs about five times that in CPU time. For the research communities this
benchmark is about, zero-shot models are already the affordable option; the question this paper answers is what
that affordability currently buys on their own data.
